\section{Introduction}
\label{sec:intro}
Non-terrestrial access can restore a path when terrestrial radio fails, or it can add delay and capacity shortfalls that push a session below its minimum-useful operating point~\cite{3gpp_tr38821,3gpp_ts22261,itur_m2160}.
Paper~III therefore treats NTN as one path among terrestrial, offline, and (named, not separately simulated here) edge/peer alternatives.
The hypothesis to reject is ``NTN always improves min-useful service versus a terrestrial baseline.''

The runnable engine is \texttt{src/ntn\_resilience/experiment.py} with policies \texttt{terrestrial\_baseline}, \texttt{static\_ntn}, \texttt{fallback}, and \texttt{adaptive}.
Sweeps cover NTN latency, capacity, and visibility; a compound power+radio failure; a configured GEO RTT contrast; and a below-min capacity probe.
Nothing in this manuscript is an operator KPI or a field satellite measurement.
